\section{Introduction and Positioning}
Can a fly-connectome spiking neural network retain collision avoidance after connection loss? The task outcome depends on the damaged network's interaction with sensing, feedback and delayed actuation. Two control laws can complete the same healthy tasks yet respond differently to matched damage. Assessing this response requires both collision outcomes and executable checks that healthy and damaged runs form the intended pairs over the declared observation horizon.

Connectome-derived models predict sensorimotor outputs\cite{shiu2024drosophila}. An reports task- and state-dependent function preservation under sparsification\cite{an2026sparse}; Jo compares circuit and virtual-arena readouts after connectome damage\cite{jo2026simulating}. Closed-loop fault evaluation is also established: DriveFI searches faults and scenarios that violate autonomous-driving safety constraints\cite{jha2019drivefi}. SpikeFI, SpikingJET and enpheeph provide configurable SNN fault campaigns\cite{spyrou2024spikefi,gogebakan2024spikingjet,colucci2022enpheeph}. These precedents motivate measuring the complete damage-to-task chain, rather than claiming a new fault-injection principle.

CRN synchronization\cite{glasserman1992crn}, hybrid-trajectory monitoring\cite{donze2010breach} and task/resource benchmarking\cite{yik2023neurobench,dangelo2026embodied} supply established tools. Cross-layer NoC analysis connects hardware routing defects to SNN behavior\cite{chessa2026crosslayer}; our functional row deletions supply no comparable physical fault mapping. Here, the evaluated increment is system-level quantification together with executable checks of the evidence supporting it.

We quantify paired task loss under two frozen feedback laws and separately recalibrated, same-prior receptive-field (RF) realizations. We contribute an evidence interface that checks campaign invariants and delegates native input/plant semantics to independent replay, with tested reuse on a delayed tank. We also evaluate frozen neural warnings against later collision, retaining misses, false alarms and uncertain transfer. Together, these contributions connect conditional task outcomes to testable pairing, timing and exposure assumptions.
